%% sections/01-introduction.tex

\section{Introduction}
\label{sec:introduction}

Resource management mechanics --- systems that constrain player actions through
finite, consumable resources --- are fundamental to game design~\cite{juul2005half,
fullerton2008}. Their effects on player behavior have been studied primarily in
terms of \textit{tactical decision-making}: resource pressure forces players to
make harder tactical choices, increasing difficulty and engagement with the
game's mechanical systems~\cite{schell2008art}.

Less studied is the effect of resource exhaustion on \textit{character-driven
decision-making} --- choices that come from character identity, relationships,
and values rather than from tactical calculation. The conventional assumption is
that resource pressure crowds out character-driven choice: when survival is at
stake, tactical calculation dominates~\cite{nacke2011brainhex}. We test this
assumption against data from the Marathon D\&D Workshop and find the opposite:
resource exhaustion in long-form campaigns amplifies character-driven decision-making
rather than suppressing it.

\subsection{The Core Puzzle}

Campaign 4 of the Marathon workshop (the-relief, Thorngate Watch) featured a
7-session no-long-rest mechanic: party HP carried forward from session to session
with only partial recovery, and the Supply Die degraded one category per session
with combat, from d12 to exhausted by session 5. By conventional game design
intuition, this level of resource pressure should produce sessions focused on
survival mechanics with reduced capacity for the emotional, character-driven
moments the rubric scores as Character Agency.

Instead, Character Agency scores rose monotonically from 7 (session 1) to 10
(sessions 4--7), reaching the maximum for the last four sessions --- including
the sessions with the most resource pressure (Supply Die at d4 and exhausted).

Campaign 6 (the-unit, Military Unit) used a Unit Cohesion Die tracking the
number of available soldiers, degrading when units were lost or orders defied.
By session 7, the Cohesion Die was at d8. Character Agency followed the same
monotone-rising pattern as Campaign 4: 7$\to$8$\to$9$\to$10$\to$10$\to$10$\to$10.

\subsection{Hypothesis}

We propose that resource exhaustion amplifies Character Agency when two conditions hold:

\begin{enumerate}
  \item \textbf{Legibility:} The party can directly observe the resource's depletion
    and its cost. A resource that degrades invisibly or abstractly does not produce
    the amplification effect.
  \item \textbf{Emotional architecture:} The campaign spine provides a reason to care
    about the resource beyond tactical survival. When the resource represents people,
    relationships, or commitments (not just power), its depletion triggers character-driven
    decision-making because character identity is implicated in resource management.
\end{enumerate}

The Supply Die (tracking garrison healing materials) and the Cohesion Die (tracking
available unit members) are both legible and emotionally architectured. The Supply Die
is legible (party can see d12$\to$d10$\to$d8$\to$d6$\to$d4$\to$exhausted) and emotionally
architectured (the resources represent people who might die without healing). The
Cohesion Die is legible and more strongly architectured (each die degradation represents
unit members who are gone).

\subsection{Contributions}

\begin{enumerate}
  \item \textbf{Empirical data:} Character Agency scores for all 49 sessions across 7
    campaigns, with resource-mechanic status as covariate.
  \item \textbf{Amplification mechanism:} An account of why resource exhaustion amplifies
    character-driven decision-making in the presence of emotional architecture.
  \item \textbf{Unit Cohesion Die advantage:} Evidence that tracking other people (rather
    than self-HP) produces stronger amplification in sessions 5--7.
  \item \textbf{Design guidance:} Implications for resource mechanic design in long-form
    narrative games.
\end{enumerate}
